\section{The Stopping Rule in Plain Terms}
\label{sec:model}

The rule we test has one moving part: a \emph{reservation price}. At every step the
agent computes the price it would expect to end up paying if it kept searching, and
buys whatever is in front of it if that offer is at least as good.

\begin{quote}
\textbf{Buy if the offer in hand is no worse than what continuing is worth.}
\end{quote}

Everything else is bookkeeping about what "continuing is worth" means when search can
fail.

\subsection{What makes continuing worth less than it looks}

Three forces push the reservation price up, making the agent more willing to accept:

\begin{description}
  \item[Fewer merchants left.] With nobody else to ask, the reservation price is just
    the offer in hand, and the agent accepts. This is why the last reachable query
    always ends in a purchase if a purchase is possible.
  \item[Risk that the next offer is unavailable.] A merchant that may be out of stock
    contributes its price only part of the time and the fallback value the rest. In
    our corpus this risk is calibrated from the measured delisting rate.
  \item[A budget that will run out.] Resources are part of the state, not a penalty
    term. A merchant that cannot be afforded is not a worse option; it is not an
    option, and it drops out of the calculation entirely.
\end{description}

The one force pushing the reservation price down is genuine price dispersion. If
another merchant is likely to be meaningfully cheaper, continuing is worth more than
the offer in hand and the agent keeps looking. This is the entire mechanism, and it
explains the central empirical finding in advance: when merchants charge the same
price, there is nothing for the rule to exploit.

\subsection{Why the budget cannot be traded away}

Time, tokens, API calls, and spend are hard limits, not costs to be balanced against
price. The agent reserves the full permit for a query \emph{before} dispatching it,
so a call can never begin unless it is affordable in the worst case. When a call is
cancelled or truncated and its true usage is unknown, it is charged at the full
permit rather than at an optimistic estimate.

The consequence is that budget overruns are impossible by construction rather than
unlikely in practice, and we treat any overrun as a correctness bug. Across every
episode and every policy replayed in this paper, there were
\ResultHardBudgetViolations{} violations.

\subsection{A rule you can compute by hand}
\label{sec:rule}

The dynamic program needs a forecast for every merchant. In practice an engineer
knows three cruder things, and those are enough.

\paragraph{Step one: the budget is one number, not four.}
Resources are consumed in fixed amounts per query, so a budget over time, tokens,
calls, and spend collapses to a single integer: how many further queries are
affordable. It is a minimum, not a sum, because whichever resource runs out first
is the one that binds.

\begin{equation}
  k = \min_{r} \left\lfloor \frac{B_r}{c_r} \right\rfloor
  \label{eq:horizon}
\end{equation}

This is a modelling convention rather than a result. It assumes deterministic
per-query costs, no partial queries, no substitution between resources, and a
budget that is never replenished.

\paragraph{Step two: accept below a fraction of the price range.}
With $k$ further queries affordable and prices expected to fall in
$[p_{\min}, p_{\max}]$, accept the offer in hand when

\begin{equation}
  p \;\le\; p_{\min} + u_k\,(p_{\max} - p_{\min}),
  \qquad u_0 = 1, \quad u_{k+1} = u_k - \tfrac{1}{2}u_k^2 .
  \label{eq:closedform}
\end{equation}

The sequence falls quickly at first and then flattens: $u_1 = \tfrac12$,
$u_2 = \tfrac38$, $u_3 = \tfrac{39}{128}$, and $u_k \sim 2/k$ thereafter. So with
one fallback left you accept the midpoint of the range; with six, only the cheapest
fifth of it.

Note that \eqref{eq:closedform} contains no per-query cost term. Cost does not enter
the objective; it enters only through \eqref{eq:horizon}, which converts the budget
into a horizon. A model charging inspection cost in the objective would instead give
$R_k = \mathbb{E}[\min(P, R_{k-1})] + c$.

\paragraph{This is a known result.}
Equation~\eqref{eq:closedform} is not new. Substituting $\mu_k = 1 - u_k$ turns it
into the classical Cayley--Moser recursion
$\mu_n = \mathbb{E}[\max(X, \mu_{n-1})] = \tfrac12(1 + \mu_{n-1}^2)$ for uniform
offers, a problem posed by Cayley in 1875 and solved by Moser in 1956, treated at
length by \citet{gilbert1966recognizing} and surveyed
by~\citet{ferguson1989secretary}. The asymptotic $u_k \sim 2/k$ is the leading term
of a published expansion~\citep{finch2024quadratic,entwistle2022asymptotics}, and the
recursion remains under active study~\citep{katriel2025cayley}. The declining
finite-horizon reservation value is standard in search
theory~\citep{kohn1974theory,lippman1976job}.

We claim only the instantiation: mapping a multi-resource query budget onto the
Cayley--Moser horizon, checking the sequence against our solver, and measuring what
it is worth on merchant data. Appendix~\ref{app:proofs} gives the derivation.

\subsection{Why the classical secretary rule is the wrong tool here}

The most widely known stopping result says to inspect about $37\%$ of the candidates
and then take the next record. It is tempting to reach for it, and it answers a
different question: it maximizes the \emph{probability of selecting the single best}
candidate using only ordinal comparisons, with no knowledge of price levels.

Our objective is the expected price paid, and our agent observes actual prices. Under
a cardinal objective the optimal cutoff is not $n/e$~\citep{bearden2006secretary}, and
the gap between rank-only and full-information sequential search is
established~\citep{palley2014sequential}. Applying the $37\%$ rule here is therefore a
category error rather than a fair contest. Section~\ref{sec:results} reports what that
error costs, because the rule is popular enough that the size of the penalty is worth
knowing.

\subsection{Is the rule actually optimal?}

For instances small enough to check exhaustively, yes. We enumerate every
deterministic accept-or-continue policy and compare the best achievable expected loss
against what the solver returns, in exact rational arithmetic with no floating-point
tolerance. Across \ResultSolverCasesVerified{} registered instances, spanning ties,
stockouts, a binding price cap, and a failure penalty large enough to dominate, the
two agree exactly (claim \texttt{SOLVER-AGREEMENT-001}).

This is a check on the implementation, not a proof of general optimality, and
agreement on small fixed-order instances says nothing about adaptive merchant
routing. The recurrence itself is stated in Appendix~\ref{app:proofs}, where readers
who want the backward induction can find it. Nothing in
Sections~\ref{sec:results}--\ref{sec:implications} depends on having read it.
